\section{Representations and bounded parent dictionaries}
\label{sec:representation}

The second trial uses shared index sets to give a solution many
representations. It must generate those representations without storing
all of them. We first specify the shared sets and the leaf lists,
then the modular filters and a finite schedule of bounded dictionary
tasks. The last part of this section proves deterministic bounds for
that schedule. Lemma~\ref{lem:representation} will show that the filters
retain a representation; Section~\ref{sec:resources} will bound the
work needed to bring it to a completed check.

The trial runs alongside the compression trial in every round. It does
not test the mean deficiency $\mu$ or decide which trial should apply.
The representation method and modular filtering have antecedents in
Howgrave-Graham and Joux~\cite{HJ10}; the multilevel worst-case setting
and its weighted disjointness step are developed by Nederlof and
W\k{e}grzycki~\cite{NW21}. The construction below specifies the shared
indices, exact integer parameters, and bounded recomputation needed
for the present resource bounds.

\subsection{Shared mixers and leaf lists}\label{alg:layout}

This trial gives each parent two children and each child two leaves.
The leaves are stored; children and parents are generated from them.
We first describe their index domains and prescribed cardinalities,
then balance the unconstrained background parts so that every leaf is
small enough to store. The shared blocks, which we call \emph{mixers}, create multiple ways to
represent a solution and require explicit disjointness tests.

Select three disjoint size-$m$ sets uniformly from a random permutation.
For each, compute its number of distinct actual subset sums exactly: store
the subset-sum lists of two size-$s$ halves and count distinct weights in
their actual-weight sorted pair stream.  Name a mixer with the largest
such number $M_0$, and name the other two $M_L,M_R$ in a deterministic
order, breaking ties deterministically.  Put
\[
 N_i=N(M_i),\qquad d_i=\defect(M_i),\qquad i\in\{0,L,R\}.
\]
Continue only if
\begin{equation}\label{eq:spread}
  \max_iN_i\le (\min_iN_i)2^{\lfloor .0005n\rfloor}.
\end{equation}
Thus $d_0\le d_L,d_R$ and the deficiency spread is at most $.0005n$.

Choose a uniform subset $F$ of all input indices, negate its weights, and
replace the target by
\begin{equation}\label{alg:sign-flip}
  t'=t-\sum_{f\in F}a_f,
\end{equation}
where the sum uses the original weights.  In the rest of this trial all
weights are signed after this change.  The signed weight of any subset
$I'$ is the original weight of $I'\mathbin{\triangle}F$ minus the
original weight of $F$.  Thus the transformed instance is equivalent to
the original, and Lemma~\ref{lem:deficiency} makes all $N_i,d_i$
unchanged.

Independently guess $k_i$ uniformly from
$\{\lceil m/6\rceil,\ldots,m/4\}$, for $i=0,L,R$, and put
$\lambda_i=k_i/m$.  At the top join, the left and right masks on $M_0$
will have sizes $k_0$ and $s-k_0$.  Within parent $j\in\{L,R\}$, its
two children will have masks on $M_j$ of sizes $k_j$ and $s-k_j$.
Disjointness of these masks is checked at the corresponding join.  The
left parent never uses $M_R$, and the right parent never uses $M_L$.

Independently choose a uniform labeled partition
\[
 M_0=T_1\mathbin{\dot\cup}T_2\mathbin{\dot\cup}T_3\mathbin{\dot\cup}T_4,
 \qquad |T_a|=g.
\]
Choose integers $c_1,\ldots,c_4$ as equally
as possible with sum $k_0$, assigning the remainder in a fixed order.
The left count in $T_a$ is $c_a$ and the right count is $g/2-c_a$.
In each parent, child 1 uses $T_1,T_2$ and child 2 uses $T_3,T_4$,
one block in each of its two leaves.  In particular, within a parent the
top-mixer indices in its children are disjoint.

Independently for $j=L,R$, partition $M_j$ uniformly into two labeled
blocks $Q_{j1},Q_{j2}$ of size $s$.  Prescribe child-1 counts
$a_{j1},a_{j2}$ as equally as possible with sum $k_j$.  The corresponding
child-2 counts are $s/2-a_{j1},s/2-a_{j2}$.  Both children use the same
two blocks of $M_j$, one in each leaf; their possible overlap is exactly
what the join of the two children into their parent checks.
Figure~\ref{tab:leaf-layout} shows both parents and all eight leaves.
Each leaf also receives its own disjoint background part, defined next. A chosen subset of a top block is recorded
in the $M_0$ mask, and a chosen subset of a local block is recorded in the
$M_j$ mask.

\begin{figure}[htbp]
\centering
\small
\setlength{\tabcolsep}{5pt}
\begin{tabular}{@{}cccc@{}}
\multicolumn{4}{c}{Final join: complementary weights and disjoint masks on $M_0$}\\[2pt]
\multicolumn{2}{c}{$\nearrow$}&\multicolumn{2}{c}{$\nwarrow$}\\[-2pt]
\multicolumn{2}{c}{\fbox{Left parent}}&
\multicolumn{2}{c}{\fbox{Right parent}}\\[3pt]
\multicolumn{2}{c}{check disjointness on $M_L$}&
\multicolumn{2}{c}{check disjointness on $M_R$}\\[1pt]
$\nearrow$&$\nwarrow$&$\nearrow$&$\nwarrow$\\[-2pt]
Left child 1&Left child 2&Right child 1&Right child 2\\[2pt]
$\uparrow$&$\uparrow$&$\uparrow$&$\uparrow$\\[-2pt]
\fbox{\begin{tabular}{@{}c@{}}$T_1,\ Q_{L1}$\\[-2pt]\scriptsize own background\end{tabular}}&
\fbox{\begin{tabular}{@{}c@{}}$T_3,\ Q_{L1}$\\[-2pt]\scriptsize own background\end{tabular}}&
\fbox{\begin{tabular}{@{}c@{}}$T_1,\ Q_{R1}$\\[-2pt]\scriptsize own background\end{tabular}}&
\fbox{\begin{tabular}{@{}c@{}}$T_3,\ Q_{R1}$\\[-2pt]\scriptsize own background\end{tabular}}\\[2pt]
$+$&$+$&$+$&$+$\\[2pt]
\fbox{\begin{tabular}{@{}c@{}}$T_2,\ Q_{L2}$\\[-2pt]\scriptsize own background\end{tabular}}&
\fbox{\begin{tabular}{@{}c@{}}$T_4,\ Q_{L2}$\\[-2pt]\scriptsize own background\end{tabular}}&
\fbox{\begin{tabular}{@{}c@{}}$T_2,\ Q_{R2}$\\[-2pt]\scriptsize own background\end{tabular}}&
\fbox{\begin{tabular}{@{}c@{}}$T_4,\ Q_{R2}$\\[-2pt]\scriptsize own background\end{tabular}}
\end{tabular}
\caption{The eight stored leaf lists and their joins, read from bottom
 to top. The two boxed leaves in each column generate one child by
 pair-sum streaming. Each $T_a$ occurs once on each parent side;
 each $Q_{jr}$ occurs in both children of side $j$. The eight
 background parts are pairwise disjoint and avoid all three mixers.
 Local joins reject overlap on $M_L$ or $M_R$; the final join rejects
 overlap on $M_0$. Boxes label allowed domains; repeated labels mark shared indices.}
\label{tab:leaf-layout}
\end{figure}

\subsubsection*{Balancing the stored lists}

There remain $r=n-3m$ background indices, whose selections have no
cardinality constraints. The prescribed mixer counts may make one
parent's lists larger than the other's, or one child's lists larger
than its sibling's. We compensate by assigning more background indices
to the smaller side, first between parents and then between children.
The program makes these allocations with integer binomial logarithms:
\[
\begin{aligned}
 b_{0L}&=\left\lfloor\log \binom{m}{k_0}\right\rfloor,&
 b_{0R}&=\left\lfloor\log \binom{m}{s-k_0}\right\rfloor,\\
 b_{j1}&=\left\lfloor\log \binom{m}{k_j}\right\rfloor,&
 b_{j2}&=\left\lfloor\log \binom{m}{s-k_j}\right\rfloor
 \quad(j=L,R).
\end{aligned}
\]
Set
\[
\begin{aligned}
 r_L&=\left\lfloor\frac{r+b_{0R}-b_{0L}}2\right\rfloor,
 &r_R&=r-r_L,\\
 r_{j1}&=\left\lfloor\frac{r_j+b_{j2}-b_{j1}}2\right\rfloor,
 &r_{j2}&=r_j-r_{j1}.
\end{aligned}
\]
Thus $r_L+r_R=r$ and $r_{j1}+r_{j2}=r_j$ exactly. Split each
child's background count as equally as possible between its two leaves,
using any fixed order on the background indices.

To verify feasibility and list sizes, write
\[
\begin{aligned}
 H_{0L}&=h(\lambda_0),& H_{0R}&=h(1/2-\lambda_0),&
 \overline H_0&=(H_{0L}+H_{0R})/2,\\
 H_{j1}&=h(\lambda_j),& H_{j2}&=h(1/2-\lambda_j),&
 \overline H_j&=(H_{j1}+H_{j2})/2\quad(j=L,R).
\end{aligned}
\]
Concavity gives $\overline H_0,\overline H_L,\overline H_R\le h_*$.
Lemma~\ref{lem:entropy} translates the integer assignments into
\begin{equation}\label{eq:allocation}
\begin{split}
 r_j&=r/2+(\overline H_0-H_{0j})m+O(\log n),\\
 r_{ji}&=r_j/2+(\overline H_j-H_{ji})m+O(\log n).
\end{split}
\end{equation}
These identities express the cancellation we need: the larger a
prescribed mixer family is, the fewer unconstrained background indices
its leaf receives.

The assignments are feasible for sufficiently large $n$.  In every
entropy pair the arguments lie in $[1/6,1/3]$, so its difference in
absolute value is at most $1-h(1/6)<3/8$.  Consequently
\[
 r_j\ge .14n-\frac{3m}{16}-O(\log n)
       \ge .095n-O(\log n),
 \qquad
 r_{ji}\ge \frac{.095n-.09n}{2}-O(\log n)>0.
\]
All prescribed mixer counts are feasible as well: they are rounded
versions of fractions in $[1/6,1/3]$ of their block sizes.  The program
may simply fail the trial if any integer feasibility test fails.

A leaf list has one occurrence for each choice with its prescribed
counts in its top and local mixer blocks and an arbitrary subset of its
background part.  An occurrence retains its actual weight and mixer
masks.  Enumerate the fixed-cardinality subsets directly, for example in
lexicographic order on index combinations, rather than scanning all
unconstrained masks.  This takes polynomial overhead per generated item.
Raw counts retain occurrence multiplicity. A raw child pair chooses
one occurrence from each of its two leaves, and a raw four-leaf tuple
chooses one from each of a parent's four leaves. Its weight is the sum
of the occurrence weights, even when local-mixer masks overlap.
The balanced background allocation makes the four leaf sizes on each
side comparable. Define their common exponential scale by
\begin{equation}\label{eq:beta}
 \beta_j=r/4+\overline H_0m/2+\overline H_jm
 \le n/4+(-.75+1.5h_*)m.
\end{equation}
Every leaf on side $j$ and each two-leaf child satisfy
\begin{equation}\label{alg:leaf-size}
 \#\{\text{leaf items}\}\le u^{10}2^{\beta_j/2},
 \qquad
 \#\{\text{raw child pairs}\}\le u^{20}2^{\beta_j}.
\end{equation}
Indeed, before substituting~\eqref{eq:allocation}, the logarithm of the
leaf size is at most
\[
 H_{0j}m/4+H_{ji}m/2+r_{ji}/2+O(1).
\]
The binomial logarithm errors in each background allocation are at most
$\log(n+1)+1$, and a rounded prescribed block count changes its entropy
upper exponent by $O(1)$: for large $n$ the relevant fractions lie in
$[1/8,3/8]$, where $|h'|<3$.  More explicitly, the error in $r_j$ from
its ideal value is at most $3\log u$, and that in $r_{ji}$ after both
ideal substitutions is at most $5\log u$.  Substitution gives
$\beta_j/2+O(\log n)$ with ample room in the factor $u^{10}$ for large
$n$.  Finally $h_*<.8114$ and $m\le .24n$ imply
\[
 \beta_j/2\le .181052n.
\]
Since $\log u\le n/10^9$, every raw leaf list has fewer than
$2^{.185n}$ items for sufficiently large main-branch inputs.

\subsection{Nested modular filters}

The first filtering level constrains each child's weight; the second
constrains the sum of the two children. Their moduli are nested so that
once a parent and its first child pass, the required congruence for its
second child follows. An additional independent prime $q$ subdivides
the parent residue classes into buckets, following the additional-prime
filtering approach of Belova, Chukhin, Kulikov, and
Mihajlin~\cite[full version, Sections~1.2.3 and~4.1]{BCKM24}.
The low filters reduce raw
generation; the extra prime controls the distinct entries in the
solution's bucket.

For analysis define the real quantities
\begin{equation}\label{eq:filter-scales}
\begin{gathered}
 \tau=.0078n,\qquad f_i=\min(\tau,d_i/2),\qquad L_i=s-f_i,\\
 f_{\max}=\max_i f_i,\qquad \xi=n/4+.032m-f_{\max}.
\end{gathered}
\end{equation}
The program uses the integer parameters
\[
 \ell_i=\lfloor L_i\rfloor
 =\max\left\{\lfloor s-\tau\rfloor,
              \left\lfloor\frac{\log N_i}{2}\right\rfloor\right\},
\]
computed by bit lengths; flooring $\log N_i$ first gives the same second
term.  Since $d_0$ is smallest, $L_0$ is largest.  Sort the three
$\ell_i$ in nondecreasing order.  For each increment from the preceding
value, taking the initial preceding value to be zero, independently
sample a prime with that increment as parameter if the increment is at
least four; otherwise use the factor one.  At each sorted index let
$p_i$ be the cumulative product.  Tied parameters give the same product.
In particular,
\[
 p_L\mid p_0,\qquad p_R\mid p_0.
\]

For every fixed nonzero relevant actual-weight difference $\Delta$,
conditional on these log parameters, the construction satisfies
\begin{equation}\label{eq:filter-bounds}
 \frac{2^{L_i}}{u^5}\le p_i\le 2^{L_i},
 \qquad
 \Pr[p_i\mid\Delta]\le u^{10}2^{-L_i}
 \quad (i=0,L,R).
\end{equation}
For the size estimate there are at most three sampled factors; the
discarded increments total at most nine.  The prime lower ranges and
the floor therefore lose a factor of at most $2^{10}(4n)^3$ relative to
$2^{L_i}$, which is at most $u^5$ for large $n$.  For divisibility,
$\Delta$ has at most $2u$ distinct prime divisors.  A factor with
parameter $\ell$ divides it with probability at most
$8u^2 2^{-\ell}$, by Lemma~\ref{lem:prime}.  All sampled factors are
independent.  Divisibility by their product requires divisibility by
each sampled factor, so multiplication of these bounds, with the same
discarded-increment and floor losses, proves~\eqref{eq:filter-bounds}.
This argument does not require the sampled primes to be distinct.

Independently sample one more prime $q$, with parameter
\[
 \ell_q=\left\lfloor n/4+.032m-s+\min_i\ell_i\right\rfloor.
\]
Then $\xi-2<\ell_q\le \xi$ and $\xi\le .25768n$; in particular $\ell_q$ is
positive and linear in $n$.  The same prime estimates give
\begin{equation}\label{alg:q-bounds}
 2^\xi/u^2\le q\le 2^\xi,
 \qquad \Pr[q\mid\Delta]\le u^3 2^{-\xi}
\end{equation}
for each fixed nonzero relevant difference.  Write $P=p_0q$.  Choose
independent uniform residues $\rho_0\pmod {p_0}$ and
$\rho_j\pmod {p_j}$ for $j=L,R$, independently also of $q$.  Their
desired parent low bases are
\[
 w_L=\rho_0,\qquad w_R=t'-\rho_0.
\]
For side $j$, require child 1 to have weight $\rho_j\pmod {p_j}$ and
child 2 to have weight $w_j-\rho_j\pmod {p_j}$.  The independent draw of
$q$ may be deferred until after all these choices when analyzing them.

\subsection{Parent dictionaries and prefix recomputation}
\label{alg:parents}

A parent dictionary stores records $(x,D)$: the exact signed integer
weight $x$ and the indexed top mask $D\subseteq M_0$. Its key is both
fields, so equal weights alone do not cause deduplication. A parent
task consists of a side, a target residue modulo $P$, and a possibly
empty prefix of the $M_0$ mask in a fixed index order. Left tasks use
the empty prefix; right tasks may use a longer prefix. The task regenerates its four leaves, joins its two child
streams, and either completes its dictionary or reports overflow.

The final check, specified in Lemma~\ref{lem:separator}, tests whether two
parent arrays contain disjoint top masks whose actual weights add to
$t'$.  Its four random separating families are sampled once before the
following loop, with the distribution specified in
Section~\ref{sec:separator}. Conditional on the layout, they use fresh
randomness independent of every parent-generation choice, including
the moduli, low residues, and extra prime, and are reused for all checks.
Their values do not affect which parent tasks are generated or how
those tasks process entries; a check can only stop the trial with YES.

For $z=0,\ldots,q-1$, set
\[
 v_L=\rho_0+p_0z,
 \qquad v_R=(t'-v_L)\bmod P,
\]
with normalized residues.  On each side these range exactly once over
the $P$-classes with the required parent low base.  Since $p_j\mid p_0$,
the child low requirements do not change with $z$.

First perform the left task.  Its two child streams are obtained from
their leaf pairs by Lemma~\ref{lem:streams}, with modulus $P$ and low
modulus $p_L$.  Stream the first child by its weight modulo $P$, and
the second by $v_L$ minus its weight modulo $P$.  Join equal keys.
For each matched pair, check that its two $M_L$ masks are disjoint.
If so, retain just the actual total weight and the combined $M_0$ mask,
deduplicating these pairs in a dictionary of capacity $C$.  If a new
distinct pair would overflow the dictionary, abandon this $z$.  Also
skip $z$ if the completed dictionary is empty.

Otherwise retain the completed left array and perform the analogous
right task with residue $v_R$, initially without an additional top-mask
restriction.  If its completed dictionary has between one and $C$
distinct entries, apply the final check to that right array and the
stored left array.  If it is empty, do nothing.  If it overflows,
discard the task's workspace and replace the task by two right tasks
that fix the next bit of its $M_0$ mask to zero and one, respectively.
Process these tasks depth first, in a fixed order of the top-mixer bits.
An overflowing task whose prefix already specifies the entire mask is
simply dropped.  Reuse or free the workspace after completing or
discarding each task, and stop at the first YES.

In every prefix task, remove incompatible items from each
leaf array \emph{before} generating its streams.  Test in that leaf
exactly the prescribed prefix bits belonging to its own top block.
Regenerating and filtering its raw leaf list from scratch is permitted.
As each top block occurs in exactly one leaf of a parent, no raw
four-leaf tuple with an incompatible top mask is generated in that
task.  This remains true even for tuples rejected later because their
local-mixer masks overlap.  It is this early restriction that permits
the later charging of repeated four-leaf tuple processing by depth
rather than by number of tasks.

\begin{lemma}\label{alg:validity}
Every retained parent entry has a subset witness on its parent's
indices.  Deduplicating entries by actual weight and top mask preserves
the existence of every valid final match.  Every YES returned by either
trial is correct.
\end{lemma}

\begin{proof}
Within a parent the children share only $M_j$, and their disjointness
there has been checked.  Their top blocks and background parts are
disjoint.  Thus a retained entry has a subset witness.  The only
indices shared by the left and right parents are in $M_0$.  Once their
top masks are fixed, the choice among witnesses of a given actual
weight has no effect on compatibility with the other parent.  Hence
deduplication loses no valid match.  A final check returning YES has
disjoint top masks and actual total weight $t'$, so its two witnesses
give a signed solution, and~\eqref{alg:sign-flip} gives an original
solution.  The compressed-block trial was already seen to have the
same one-sided correctness property.
\end{proof}

\subsection{Three bounds for prefix recomputation}
\label{sec:prefix-accounting}

The schedule has two different sources of work. Starting a task
regenerates leaves and initializes streams; processing its matches
examines raw four-leaf occurrences. A small dictionary can have many
raw representations, including invalid local overlaps. We therefore
use distinct entries to count task starts and raw occurrences to count
match processing. A third count measures how often retained arrays
are passed to the final check.

\paragraph{What the trial counts.}
Fix the mixers, signs, cardinality guesses, layout, low moduli, and
low residues, without fixing the extra prime $q$ or the separating
families. The following three finite counts are already determined at this stage:
\begin{itemize}
\item $G$ is the sum, over the four children, of the numbers of raw
leaf pairs passing that child's low filter, before any mask-prefix
restriction.
\item $\mathcal T_j$ contains all raw four-leaf occurrences on side
$j$ passing the parent low filter modulo $p_0$ and the first-child
filter modulo $p_j$. The second-child filter follows from
$p_j\mid p_0$. Put $J_j=|\mathcal T_j|$, retaining every repetition
and every local overlap.
\item $\mathcal U_j$ contains all distinct valid
$(\text{actual weight},\text{top mask})$ entries satisfying the
parent low filter, ignoring child filters. Here validity means that
the entry has a subset witness on side $j$'s indices with the
prescribed top-block counts. This ambient set may contain entries
not generated by the child lists. Put $U_j=|\mathcal U_j|$.
\end{itemize}
After $q$ is chosen, assign each occurrence and entry its unique
bucket using the enumerated $P=p_0q$ classes. On the right the bucket
label is the value of $z$ whose class is
$(t'-\rho_0-p_0z)\bmod P$. These labels depend on $q$, but the three
counts do not. Equal entry keys have equal actual weights and hence
the same bucket. No coprimality of $p_0$ and $q$ is needed.

The next lemma separates these deterministic charges from all
probability estimates. The schedule is the one specified in
Section~\ref{alg:parents}: left overflow skips a bucket, whereas right
overflow refines a mask prefix. Only for counting, continue this
schedule after any YES answer and ignore the operation cap. This
finite hypothetical schedule need not be stored.

\begin{lemma}[Prefix accounting]\label{lem:prefix}
Fix positive integers $C,q$ and a nonnegative integer $m$. On each
side $j\in\{L,R\}$, let $\mathcal T_j$ be a finite collection of
raw tuple occurrences. Each occurrence has a bucket in
$\{0,\ldots,q-1\}$ and a top mask in $\{0,1\}^m$. A valid occurrence
produces an entry consisting of its exact weight and top mask.
Assume that occurrences producing the same entry have the same
bucket. Let $\mathcal U_j$ be a finite ambient set containing all
distinct valid entries, with bucket and mask labels extending those
of the generated entries. Write
\[
 J_j=|\mathcal T_j|,\qquad U_j=|\mathcal U_j|.
\]
A task for bucket $z$ and prefix $\sigma$ generates exactly the raw
occurrences with these labels, unless it stops early, and reports
overflow at the $(C+1)$st distinct valid entry. Suppose a task with
$Z$ potential occurrences costs at most $H+KZ$, even if stopped
partway, for fixed nonnegative $H,K$.

Apply the left-skip and right-prefix schedule of
Section~\ref{alg:parents} to these data, ignoring early YES answers and
the operation cap. Then its number $T$ of parent-generation tasks,
its total parent-generation work $W_{\mathrm{parent}}$, and the total
input-array length $N_{\mathrm{check}}$ of its final checks satisfy
\begin{align}
 T&\le 2q+2(m+1)U_R/C,\label{eq:prefix-tasks}\\
 W_{\mathrm{parent}}
 &\le H\bigl(2q+2(m+1)U_R/C\bigr)
       +K\bigl(J_L+(m+1)J_R\bigr),\label{eq:prefix-work}\\
 N_{\mathrm{check}}&\le U_L+(2m+3)U_R.\label{eq:prefix-checks}
\end{align}
Every reuse of a left array counts separately in
$N_{\mathrm{check}}$.
\end{lemma}

\begin{proof}
\emph{Task starts.}
At a fixed prefix depth, different right tasks have disjoint sets of
possible entries: different buckets are disjoint, and different
prefixes of the same length are disjoint within a bucket. Each
overflow witnesses more than $C$ entries of $\mathcal U_R$.
There are consequently at most $U_R/C$ overflow nodes at that depth,
and at most $(m+1)U_R/C$ over all depths. There are at most $q$ left
tasks and $q$ right root tasks. Every other task is one of at most two
children of an overflow node. This proves~\eqref{eq:prefix-tasks},
including the possibility that an overflowing full-mask node is
dropped without children.

\emph{Raw occurrences.}
Each left occurrence has one bucket and can appear in at most one
left task. A right occurrence has one bucket and one prefix at each
depth, so it is a potential occurrence in at most $m+1$ right tasks.
This counts occurrence identities, rather than their possibly equal
weights or masks. It also counts invalid occurrences: their top masks
are defined before the local overlap test, and incompatible prefix
choices are excluded before the join. Thus the sum of all potential
occurrence counts is at most $J_L+(m+1)J_R$. Summing the assumed
stopped-task costs gives~\eqref{eq:prefix-work}. Only the overhead
$H$ is multiplied by the number of tasks.

\emph{Final-check inputs.}
Completed right arrays contain disjoint sets of entries. Different
buckets are disjoint. Within a bucket, completed prefixes are
incomparable, since a completed task has no descendants, so their
mask classes are disjoint as well. Right arrays therefore contribute
at most $U_R$ to the total check-input length.

Left arrays used at right-root checks have total length at most $U_L$,
since each bucket supplies at most one such check. Every other check
uses a left array of length at most $C$, and there are at most
$2(m+1)U_R/C$ nonroot tasks. Their repeated left lengths total at most
$2(m+1)U_R$. Adding these three contributions gives
\eqref{eq:prefix-checks}.
\end{proof}

\paragraph{Applying the accounting bounds.}
Every raw four-leaf tuple has a well-defined top mask because
$T_1,\ldots,T_4$ partition $M_0$ and each occurs in exactly one leaf
on each side. The union is defined even if local-mixer masks overlap.
Testing a prefix in its responsible leaves therefore gives exactly
the restricted occurrence collection required by Lemma~\ref{lem:prefix}.
The dictionary removes duplicate keys only after the local
disjointness test. This verifies the combinatorial premises of the
lemma for the concrete schedule.

\paragraph{The stopped-task cost.}
For sufficiently large main-branch inputs, its cost premise holds with
\begin{equation}\label{eq:prefix-task-cost}
 H=u^8(2^{.2n}+G),\qquad K=u^8.
\end{equation}
Here is the implementation behind those constants. The four raw leaf
arrays in a task have total length
$A<4\cdot2^{.185n}$ by~\eqref{alg:leaf-size}; regenerating them,
including their masks and prefix tests, takes $O(n^2A)$ operations.
Sorting and stream initialization cost $O(nA)$, and the two child
streams together have at most $G$ occurrences. Their arrays and
mutable states occupy $O(A)$ words, which is at most $B$ for large
$n$. The tie buffer has length $B$, and Lemma~\ref{lem:streams}
charges all replay, including stopped replay, to stream advances and
potential matching pairs. Dictionary initialization can be bounded by
$O(u(C+1))$, and local mask checks and trie operations cost $O(u)$ per
reported tuple. For $Z$ potential matching tuples, all this gives
\[
 O\bigl(n^2A+B+nG+u(C+1)+uZ\bigr).
\]
Since $C\le2^{.199n}$, $B\le2^{.19n}$, and
$\log u\le n/10^9$, this is at most $H+KZ$ with
\eqref{eq:prefix-task-cost}, after enlarging the absolute main-branch
threshold. The raw term includes tuples later rejected for overlap
and repeated occurrences. A small value of $U_j$ does not bound this
term.

Thus Lemma~\ref{lem:prefix} controls the complete schedule before any
probability estimate is used. The numerical bounds for $G,J_j,U_j$
and the proof that a solution reaches a completed right prefix belong
to Section~\ref{sec:resources}. The remaining independent ingredient
is the final check on two bounded arrays, which we now construct.
